% GDZ PPN236018647,native510;HKUST511;original diary12r,service25.
% Complete entries141-143,145-146;144 cancellation unresolved.
\clearpage\SourceLeaf{510}\NoSourceHead
\begingroup\fontsize{10}{15}\selectfont\itshape
\noindent Catalogum praecedentem per fata iniqua iterum interruptum
initio anni 1812 resumimus. In mense Nov. 1811 contigerat
demonstrationem theorematis fundamentalis in doctrina aequationum
pure analyticam completam reddere; sed quum nihil
chartis servatum fuerit, pars quaedam
\tikz[baseline=(m.base)]{
\node[anchor=base west,inner sep=0pt,font=\scriptsize\itshape] (e) at (0,1.4ex) {essentialis};
\node[anchor=base west,inner sep=0pt,font=\itshape] (m) at (0,0) {memoriae};
\draw[line width=.4pt] (-.2em,-.6ex) .. controls (-.5em,1ex) and (.2em,1ex) .. (e.south east);
}
penitus exciderat. Hanc per satis longum temporis intervallum
frustra quaesitam tandem feliciter redinvenimus 1812 febr. 29

\medskip
\noindent\underline{Theoriam Attractionis Sphaeroidis Elliptici in puncta}\\
\underline{extra solidum sita prorsus novam invenimus}
\par\noindent\hfill Seeberg. 1812. Sept. 26

\medskip
\noindent\underline{Etiam partes reliquas eiusdem theoriae per methodum}\\
\underline{novam mirae simplicitatis absolvimus}\hfill 1812 Oct. 15 Gott

\medskip
\noindent Subtilissimum hoc est omnium eorum quae umquam
perfecimus. Vix itaque operae pretium est, his intermiscere
mentionem quarumdam simplificationum ad calculum
orbitarum parabolicarum pertinentium.

\medskip
\noindent Observatio per inductionem facta gravissima theoriam residuorum biquadraticorum
cum functionibus lemniscaticis elegantissime nectens. Puta si \(a+bi\) est
numerus primus \(a-1+bi\) per \(2+2i\) divisibilis, multitudo omnium solutionum
congruentiae \(1=xx+yy+xxyy\ (\mathrm{mod}\ a+bi)\), inclusis \(x=\infty,\ y=\pm i,\)
\(x=\pm i,\ y=\infty\) fit \(=(a-1)^2+bb\)
\par\noindent\hfill 1814 Jul 9
\endgroup
